\chapter{Literature Survey}

\section{Review of Recent Literature}

\subsection{Machine Learning in Agriculture}

Machine learning in agriculture was reviewed in broad terms by Liakos et al.~\cite{liakos2018ml}, who grouped the work by
application: crop management, livestock, water and soil management. Kamilaris and Prenafeta-Bold{\'u}~\cite{kamilaris2018deep}
surveyed deep learning in agriculture and found that convolutional networks dominated image-based tasks such as
disease detection and weed identification. For plant disease detection in particular, Mohanty et al.\ trained deep
networks on the PlantVillage image collection and reported high accuracy on held-out images from the same
collection~\cite{mohanty2016plantvillage,hughes2015plantvillage}, and Ferentinos~\cite{ferentinos2018deep} compared several
architectures on a larger set of leaf images. These works establish that supervised image classifiers can be accurate
when images resemble the training data (single leaves on plain backgrounds), but they output a class label only: no
explanation, no treatment advice and no link to the conditions in the field.

\subsection{Large Language Models for Agricultural Question Answering}

Silva et al.\ examined whether a general-purpose model, GPT-4, can answer agriculture examination questions~\cite{silva2023gpt4agronomist}.
Balaguer et al.\ compared retrieval-augmented generation with fine-tuning in a pipeline study that used agriculture as its
case study, and reported that adaptation to local, geography-specific knowledge improved the similarity of answers to
reference answers~\cite{balaguer2024ragvsft}. Yang et al.\ presented ShizishanGPT, an agricultural assistant that combines
a general model with a search engine, a knowledge graph, a retrieval module and an agent that can call specialised
crop models~\cite{yang2024shizishangpt}. Together these studies show that retrieval and tool use are the accepted way to
bring domain and local knowledge into a general model. They also share a limitation relevant to this project: the
knowledge comes from documents, benchmarks and public sources, not from the \emph{farmer's own live measurements}.

\subsection{Vision-Language Models for Plant Disease}

Two recent benchmarks evaluate general vision-language models on agricultural images. Ranario and Earles benchmarked open and
closed models on many agricultural classification datasets and concluded that, used zero-shot, the models trail a
supervised baseline and perform much worse with open-ended prompting than with constrained, multiple-choice
questions. They recommend using such models with constrained interfaces and not as stand-alone diagnostic
tools~\cite{ranario2025vlmag}. LeafNet introduced a large visual question-answering benchmark for plant diseases and
reported that models do well on the coarse healthy-versus-diseased decision but considerably worse when asked to identify the
pathogen or the species~\cite{nguyen2026leafnet}. These findings directly shaped the design of Pipeline B: a strict
JSON schema constrains the model, the output separates observed evidence from assumptions, the diagnosis is phrased as a
possibility, and a list of look-alike conditions is requested.

\subsection{Confidence and Uncertainty}

Neural classifiers are often over-confident, and calibration is commonly measured with the expected calibration error
and reliability diagrams~\cite{guo2017calibration}. Language models can be asked to state their confidence, and studies
show that these verbalised numbers carry some signal but are imperfect~\cite{kadavath2022know,xiong2024uncertainty}. For this
reason GAIA never presents the model's self-reported confidence as a probability: it is shown as a rough indicator with
a fixed caveat, and the evaluation scripts compute the calibration metrics so that the caveat is backed by
measurements.

\begin{table}[H]
\centering
\caption{Summary of the reviewed literature and its relevance to GAIA.}
\label{tab:litsummary}
\small
\begin{tabular}{|P{2.4cm}|P{4.0cm}|P{5.9cm}|}
\hline
\textbf{Work} & \textbf{Focus} & \textbf{Relevance to GAIA}\\\hline
Liakos et al.~\cite{liakos2018ml}; Kamilaris and Prenafeta-Bold{\'u}~\cite{kamilaris2018deep} & Reviews of machine learning and deep learning in agriculture & Establish the application area; image classifiers dominate but give labels only\\\hline
Mohanty et al.~\cite{mohanty2016plantvillage}; Ferentinos~\cite{ferentinos2018deep} & CNN disease classifiers on leaf images & PlantVillage is used in GAIA's vision benchmark script; results are an upper bound for field photos\\\hline
Silva et al.~\cite{silva2023gpt4agronomist} & GPT-4 on agriculture exams & Motivates grounding a general model in domain data\\\hline
Balaguer et al.~\cite{balaguer2024ragvsft} & RAG versus fine-tuning, agriculture case study & Supports choosing RAG with prompts over fine-tuning\\\hline
Yang et al.~\cite{yang2024shizishangpt} & Agricultural LLM with search, knowledge graph, RAG and agent & Closest system: tool-using assistant, but no live farm sensor data\\\hline
Ranario and Earles~\cite{ranario2025vlmag} & Zero-shot VLMs versus supervised models & Motivates constrained, strict-JSON prompting and honest framing\\\hline
Nguyen Quoc et al.~\cite{nguyen2026leafnet} & Plant-disease VQA benchmark & Coarse health is reliable, fine identification is not: show possibilities and differentials\\\hline
Guo et al.~\cite{guo2017calibration}; Xiong et al.~\cite{xiong2024uncertainty} & Calibration and confidence elicitation & Justify the uncalibrated-confidence caveat and the ECE and Brier metrics\\\hline
\end{tabular}
\end{table}

\section{Generative AI: Existing Approaches}

\subsection{Prompting, Retrieval, Fine-Tuning and Agents}

Four approaches are commonly used to adapt a general model to a domain. Table~\ref{tab:approaches} compares them for the
needs of this project.

\begin{table}[H]
\centering
\caption{Approaches to adapting a generative model to a domain.}
\label{tab:approaches}
\small
\begin{tabular}{|P{2.6cm}|P{3.7cm}|P{3.7cm}|P{2.4cm}|}
\hline
\textbf{Approach} & \textbf{Strengths} & \textbf{Weaknesses for GAIA} & \textbf{Used?}\\\hline
Plain prompting & Simple, cheap & Cannot know the farm's data & Only for small talk\\\hline
Retrieval-augmented generation~\cite{lewis2020rag} & Uses fresh private data; answers can cite sources & Needs good retrieval; model can still ignore context & Yes (Pipeline A)\\\hline
Fine-tuning & Learns style or domain vocabulary & Cannot absorb live sensor data; needs data and compute & No\\\hline
Agents with tools~\cite{yao2023react} & Can decide what to look up & Unbounded loops, cost, harder to test & Yes, bounded (Pipeline B)\\\hline
\end{tabular}
\end{table}

\subsection{Retrieval Techniques}

Dense retrieval embeds text into vectors and finds nearest neighbours by cosine similarity. Approximate nearest-neighbour
indexes such as Hierarchical Navigable Small World (HNSW) graphs make this fast for large collections~\cite{malkov2020hnsw},
and the pgvector extension provides them inside PostgreSQL~\cite{pgvector}. Dense retrieval captures meaning (``dry
soil'' is close to ``low soil moisture'') but can miss exact tokens such as sensor names. Lexical search has the opposite
profile. Combining the two rankings with Reciprocal Rank Fusion (RRF) is a simple, parameter-light method that was shown
to outperform several learned fusion methods~\cite{cormack2009rrf}. GAIA uses this hybrid, with the Gemini embedding model
providing the vectors~\cite{lee2025geminiembedding}.

\subsection{Multimodal Models}

Contrastive pre-training on image--text pairs~\cite{radford2021clip} and visual instruction tuning~\cite{liu2023llava} led to
models that can describe and reason about images. Commercial multimodal models such as Gemini accept an image and a
text instruction together and can return structured output~\cite{gemini2023,gemini25,geminiapi}. GAIA uses this
capability in a single call to analyse a crop photograph.

\section{Review of Related Generative AI Systems}

\begin{table}[H]
\centering
\caption{Related systems compared with GAIA, based on the cited publications. A dash means the publication does not report the feature and no claim is made.}
\label{tab:related}
\small
\begin{tabular}{|P{3.1cm}|P{1.7cm}|P{1.7cm}|P{1.8cm}|P{2.0cm}|P{1.9cm}|}
\hline
\textbf{System} & \textbf{Image diagnosis} & \textbf{Language Q\&A} & \textbf{Web or document retrieval} & \textbf{Own live sensor data} & \textbf{Explains data gaps}\\\hline
Plantix~\cite{plantix} & Yes & -- & -- & No & No\\\hline
CNN classifiers~\cite{mohanty2016plantvillage,ferentinos2018deep} & Yes (label) & No & No & No & No\\\hline
GPT-4 on exams~\cite{silva2023gpt4agronomist} & No & Yes & -- & No & No\\\hline
ShizishanGPT~\cite{yang2024shizishangpt} & Not reported & Yes & Yes & No & No\\\hline
\textbf{GAIA (this work)} & \textbf{Yes (VLM)} & \textbf{Yes} & \textbf{Yes} & \textbf{Yes (rover)} & \textbf{Yes}\\\hline
\end{tabular}
\end{table}

\section{Existing Solutions / State-of-the-Art Methods}

\begin{itemize}
\item \textbf{Consumer disease-identification applications.} Photograph in, label and advice out~\cite{plantix}. They depend
      on the farmer noticing a symptom.
\item \textbf{Supervised leaf classifiers.} High accuracy on curated datasets~\cite{mohanty2016plantvillage,ferentinos2018deep}
      but little transfer to field photographs and no explanation.
\item \textbf{Retrieval-augmented agricultural assistants.} Document-grounded answers, sometimes with tools and
      knowledge graphs~\cite{balaguer2024ragvsft,yang2024shizishangpt}.
\item \textbf{General multimodal models used zero-shot.} Flexible, but weaker than supervised models on fine-grained
      agricultural classes~\cite{ranario2025vlmag,nguyen2026leafnet}.
\end{itemize}

\section{Gap Identification}

\begin{enumerate}[label=G\arabic*.]
\item \textbf{Reactive operation.} Existing diagnosis tools wait for the farmer to notice a problem. A rover that walks the
      field and measures conditions on demand supports proactive inspection.
\item \textbf{No grounding in the farmer's own measurements.} Retrieval-based agricultural assistants answer from public
      or curated documents. None of the reviewed systems answers from the farm's own sensor log with exact statistics.
\item \textbf{Numbers from language models.} Letting a model compute averages or dates invites errors. A design in which code
      computes all numbers and the model only phrases them is rarely described.
\item \textbf{Silent handling of missing data.} Systems treat a missing reading as ``no data'' without saying whether a
      sensor was switched off, silent or disconnected. This matters once battery-saving switches exist.
\item \textbf{Overstated certainty.} Diagnosis tools commonly present a single label and a percentage. Evidence from the
      literature~\cite{guo2017calibration,nguyen2026leafnet} argues for possibility-phrased output with visible evidence and
      an uncertainty caveat.
\item \textbf{Isolated pipelines.} Chat and diagnosis are normally separate products. Sharing a knowledge store lets the
      assistant later discuss a diagnosis together with the field conditions at that time.
\end{enumerate}

\section{Common Findings from Literature}

\begin{itemize}
\item Retrieval and tool use, not model size alone, are the standard way to bring domain and local knowledge into a general
      model~\cite{lewis2020rag,gao2023ragsurvey,balaguer2024ragvsft}.
\item General vision-language models are reliable for coarse health assessment but unreliable for fine-grained
      identification, so interfaces should constrain them and present results as possibilities~\cite{ranario2025vlmag,nguyen2026leafnet}.
\item Model confidence is informative but not calibrated and should be disclosed as such~\cite{guo2017calibration,kadavath2022know,xiong2024uncertainty}.
\item Hybrid lexical and dense retrieval is robust and simple to implement~\cite{cormack2009rrf}.
\item Hallucination remains a general risk~\cite{ji2023hallucination}, so outputs should be checked by code wherever
      the answer can be verified (numbers, citation identifiers).
\end{itemize}
